\clearpage
\section*{Библиография}
\addcontentsline{toc}{section}{Библиография}
\label{sec:bibliografiya}

% Списъкът е с разстояние 1,5 реда като основния текст; размерът на шрифта е 11 pt.
% Нормативните актове са по lex.bg към 16.11.2026 г.; актовете на ЕС и решението на СЕС~--
% по EUR-Lex.
{\small

\subsection*{Литература}
\begin{biblist}[series=bib]
\item Господинов, Д. Проблеми на електронните документи и електронните подписи при
  прилагането на Регламент eIDAS от българския съд.~-- Общество и право, 2024, №~7,
  с.~19--48.
\item Димитров, Г. Право на информационните и комуникационните технологии~--
  гражданскоправни аспекти. София: Фондация „Право и интернет“, 2014.
\item Димитров, Г., С.~Ракшиева. Електронна идентификация и електронни удостоверителни
  услуги.~-- Юридически сборник, 2016, т.~23, с.~100--112.
\item Калайджиев, А. Относно същността на електронното изявление и електронния
  документ.~-- Пазар и право, 2002, №~1.
\item Калайджиев, А., Б.~Белазелков, Г.~Димитров, М.~Йорданова, Д.~Марков, В.~Станчева.
  Електронният документ и електронният подпис~-- правен режим. София: Сиела, 2004.
\item Младенов, Д. Електронни доказателства.~-- De Jure, 2022, т.~13, №~1, с.~63--74.
\item Сталев, Ж., А.~Мингова, В.~Попова, О.~Стамболиев, Р.~Иванова. Българско гражданско
  процесуално право. 10~изд. София: Сиела, 2020.
\end{biblist}

\subsection*{Нормативни актове}
\begin{biblist}[resume*=bib]
\item Граждански процесуален кодекс. Обн. ДВ, бр.~59 от 20.07.2007~г., в сила от
  01.03.2008~г.; последно изм. ДВ, бр.~92 от 06.10.2026~г.
\item Закон за електронния документ и електронните удостоверителни услуги (загл. изм.,
  ДВ, бр.~85 от 2017~г.). Обн. ДВ, бр.~34 от 06.04.2001~г., в сила от 06.10.2001~г.;
  последно изм. ДВ, бр.~58 от 23.07.2019~г.
\item Закон за задълженията и договорите. Обн. ДВ, бр.~275 от 22.11.1950~г.; последно изм.
  ДВ, бр.~35 от 27.04.2021~г.
\item Наказателен кодекс. Обн. ДВ, бр.~26 от 02.04.1968~г.; последно изм. ДВ, бр.~55 от
  16.06.2026~г.
\end{biblist}

\subsection*{Актове на Европейския съюз}
\begin{biblist}[resume*=bib]
\item Регламент (ЕС) №~910/2014 на Европейския парламент и на Съвета от 23 юли 2014~г.
  относно електронната идентификация и удостоверителните услуги при електронни
  трансакции на вътрешния пазар и за отмяна на Директива 1999/93/ЕО. ОВ L~257,
  28.08.2014~г., с.~73. Консолидиран текст към 18.10.2024~г.
\item Регламент (ЕС) 2024/1183 на Европейския парламент и на Съвета от 11 април 2024~г.
  за изменение на Регламент (ЕС) №~910/2014 по отношение на създаването на европейска
  рамка за цифрова самоличност. ОВ L, 2024/1183, 30.04.2024~г.
\item Регламент за изпълнение (ЕС) 2025/1566 на Комисията от 29 юли 2025~г. за определяне
  на правила за прилагането на Регламент (ЕС) №~910/2014 на Европейския парламент и на
  Съвета по отношение на референтните стандарти за проверка на самоличността и атрибутите
  на лице, на което ще бъде издадено квалифицирано удостоверение или квалифицирано
  електронно удостоверение за атрибути. ОВ L, 2025/1566, 30.07.2025~г.
\end{biblist}

\subsection*{Съдебна и административнонаказателна практика}
\begin{biblist}[resume*=bib]
\item Решение на СЕС от 20 октомври 2022~г. по дело C-362/21, \textit{Екофрукт},
  ECLI:EU:C:2022:815.
\item Решение №~303 от 07.10.2013~г. по гр.~д. №~3715/2013~г., ВКС, IV~г.~о.
\item Решение №~70 от 19.02.2014~г. по гр.~д. №~868/2012~г., ВКС, IV~г.~о. (цит. по
  Определение №~169 от 06.04.2017~г.).
\item Определение №~169 от 06.04.2017~г. по ч.~т.~д. №~672/2017~г., ВКС, I~т.~о.
\item Решение №~176 от 02.07.2025~г. по т.~д. №~123/2023~г., ВКС, II~т.~о.
\item Решение №~164 от 16.05.2022~г. по в.~гр.~д. №~173/2022~г., Окръжен съд~--
  Пазарджик.
\item Решение №~302 от 22.10.2014~г. по а.н.д. №~802/2014~г., Районен съд~-- Хасково.
\item Решение №~22 от 10.02.2015~г. по к.а.н.д. №~374/2014~г., Административен съд~--
  Хасково.
\item Наказателно постановление №~123 от 07.07.2014~г. на председателя на Комисията за
  регулиране на съобщенията.
\end{biblist}
}
